% Generated draft: requires mathematical and source review.
\documentclass[11pt]{article}
\usepackage[margin=25mm]{geometry}
\usepackage{fontspec}
\setmainfont{Latin Modern Roman}
\usepackage{amsmath,amssymb,mathtools,unicode-math}
\setmathfont{Latin Modern Math}
\usepackage{enumitem,needspace,xurl,hyperref}
\setlength{\parindent}{0pt}
\setlength{\parskip}{5pt}
\setlength{\emergencystretch}{3em}
\newcommand{\N}{\mathbb{N}}
\newcommand{\Z}{\mathbb{Z}}
\newcommand{\Q}{\mathbb{Q}}
\newcommand{\R}{\mathbb{R}}
\newcommand{\C}{\mathbb{C}}
\newcommand{\F}{\mathbb{F}}
\newcommand{\A}{\mathbb{A}}
\newcommand{\PP}{\mathbb{P}}
\newcommand{\E}{\mathbb{E}}
\newcommand{\eps}{\varepsilon}
\newcommand{\dd}{\,\mathrm{d}}
\begin{document}
\section*{UnsolvedMath ID: 30003741}
\textbf{Source problem code:} catalogue entry
\subsection{Definitions and mathematical statement}
The supplied question concerns a T-cell activation ordinary differential equation model, with antigen amount and receptor-unbinding rate as parameters. The record does not supply the differential equations, the state space, the admissible parameter set, or the treatment of other model parameters. Consequently, the following is a schematic formulation whose missing model data must be recovered from the cited sources before it becomes a fully specified mathematical problem.

Write the intended model schematically as
\[
\frac{\dd x}{\dd t}=F(x;p),\qquad x\in\Omega\subseteq\R^n,\qquad p\in\mathcal P,
\]
where \(x\) is the state, \(\Omega\) is its admissible domain, and \(p\) includes the antigen amount and receptor-unbinding rate. Here \(n\), \(F\), \(\Omega\), and \(\mathcal P\) are unspecified model data, not additional assumptions. A steady state at a fixed parameter value \(p\) is a point \(x_*\in\Omega\) satisfying \(F(x_*;p)=0\).

The requested existence question is whether the intended model admits more than three distinct steady states for one admissible parameter choice; equivalently, whether
\[
\exists p\in\mathcal P\;\exists x^{(1)},x^{(2)},x^{(3)},x^{(4)}\in\Omega:
\quad F(x^{(i)};p)=0\ \ (1\leq i\leq4),\qquad
x^{(i)}\neq x^{(j)}\ \ (1\leq i<j\leq4).
\]
This formulation leaves unresolved whether only the two named parameters may vary, which other parameters are fixed, and whether the model permits varying the number of sites. It also leaves unresolved whether steady states must be strictly positive, may lie on the boundary of the physical state space, or must belong to a single prescribed conservation class.

The supplied dated assessment states that three positive steady states are proved possible and that an upper bound of three is proved for the three-site case, while existence of parameter values yielding more than three steady states remains unresolved. These are supplied literature claims, not independently verified mathematical conclusions. The record does not specify the model variant or hypotheses needed to reconcile the general existence question with the stated three-site bound.
\subsection{Short English statement}
Can the intended T-cell activation ODE model have at least four distinct steady states at one admissible parameter choice involving antigen amount and receptor-unbinding rate? The equations, parameter restrictions, site count, and meaning of admissible steady states are missing. The supplied record labels the problem open, but its dated claims and intended formulation require human review.
\subsection{Sources}
\begin{itemize}
\item Local registry: \texttt{data/unsolvedmath/problems.json}, supplied record with permanent UnsolvedMath ID \(30003741\), source problem code \texttt{catalogue entry}, titled ``Higher Multistationarity in T-Cell Activation Models,'' in the Oberwolfach Reports Open Problems collection. No individual proposer is supplied; the recorded proposed year is 2018.
\item Supplied source citation: \emph{The Mathematics of Mechanobiology and Cell Signaling} (2019), Oberwolfach Reports, pp.\ 433--506, DOI 10.4171/owr/2018/8: \url{https://doi.org/10.4171/owr/2018/8}. The background describes Oberwolfach Reports as documenting mathematical workshops and their open-problem discussions.
\item Supplied literature source, labeled ``original OWR report'': \url{https://doi.org/10.4171/OWR/2018/8}.
\item Supplied literature source, labeled ``Rendall--Sontag, multiple T-cell steady states'': \url{https://doi.org/10.1098/rsos.170821}. The record supplies no fuller bibliographic citation or pinpoint reference for the model or the stated bounds.
\item Literature assessment date: \texttt{2026-08-21}. This is supplied dated curation, not independently verified. It labels the problem open and reports that three positive steady states are proved possible, that at most three are proved for the three-site case, and that existence of parameter values giving more than three remains unresolved. The background explicitly describes this as dated literature triage derived from supplied OWR curation and requests independent verification of time-sensitive status and cited claims.
\end{itemize}
\end{document}
